\documentclass[10pt,a4paper]{amsart}
\usepackage[margin=19mm]{geometry}
\usepackage{amsmath,amssymb,mathtools}
\usepackage[scr=rsfs,cal=euler]{mathalfa}
\usepackage{xfrac}
\usepackage[unicode,hidelinks,bookmarks=false]{hyperref}
\newcommand{\ie}{i.e.\ }
\newcommand{\cf}{cf.\ }
\newcommand{\resp}{respectively\ }
\newcommand{\Subsection}{\subsection}
\providecommand{\nobreakdash}{\nobreak\hbox{-}}
\setlength{\emergencystretch}{2em}
\multlinegap0pt
\usepackage{amssymb}
\usepackage[shortlabels]{enumitem}
\usepackage{xcolor}
\newtheorem{lemma}{Lemma}
\newcommand{\res}{\operatorname{res}}
\newcommand{\val}{\operatorname{val}}
\title{The multiplicative structure of the K-theoretical McKay correspondence for the Hilbert scheme of points in the complex plane}
\author{Jakub Koncki, Magdalena Zielenkiewicz}
\date{Source: arXiv:2407.14293v2 (CC BY 4.0)}
\begin{document}
\maketitle

\noindent\emph{Source and licence.} The multiplicative structure of the K-theoretical McKay correspondence for the Hilbert scheme of points in the complex plane. Version arXiv:2407.14293v2 (CC BY 4.0), distributed by arXiv under the Creative Commons Attribution 4.0 International licence, http://creativecommons.org/licenses/by/4.0/, as recorded in the arXiv licence metadata for this exact version and shown on the abstract page https://arxiv.org/abs/2407.14293v2. Full licence notice: \texttt{licenses/arxiv-2407.14293-CC-BY-4.0.txt}.\par\smallskip

\noindent\textit{Editorial scope.} Counted unit: the article's lemma lemma (source0.tex lines 974--976). The article's complete original proof of it is delivered with it (source0.tex lines 977--987, one \texttt{proof} environment of 11 lines). No mathematics is written, repaired or re-derived: the statement and the proof are the article's own lines, in the article's own notation, and no retained line is substituted or re-flowed.  No further source-local context is needed: every cross-reference of every retained line resolves inside the two delivered spans.  Nothing was omitted from any retained span, so the package carries no omission record.  No retained line contains a citation, so the wrapper carries no bibliography and no bibliography entry.  No retained line contains a figure environment, an image-inclusion command or drawing code, so no image is delivered and the package declares no asset dependency.  Editorial formatting only, in addition: an AMS \texttt{amsart} preamble that restates the article's own declarations and macros used by the retained lines, namely the environments \texttt{lemma} on separate counters, together with the standard packages those lines need.  The retained environments print in this document as Lemma 1 in order of appearance, whereas the article prints the same items with the numbering of its own full document.  The complete native source of the article as distributed in the arXiv e-print for this exact version is bundled unchanged as \texttt{sources/162579/source0.tex}.
\begin{lemma}
	Let $k$ be a field, let $V$ and $W$ be $k$-vector spaces of finite dimension, and let \hbox{$\psi:V\to W$} be a $k$-linear map. Let $\psi_q:V\otimes_k k(q)\to W\otimes_k k(q)$ be a $k(q)$-linear map which restricts to $\psi$. Suppose that $\psi$ is an isomorphism. Then $\psi_q$ is also isomorphism.
\end{lemma}

\begin{proof}
	Let $v_1,\dots,v_k$ be a basis of $V$. Then $v_1\otimes 1,\dots,v_k\otimes 1$ is a $k(q)$-basis of $V\otimes k(q)$. It is enough to prove that the vectors $\psi_q(v_1\otimes 1),\dots,\psi_q(v_k\otimes 1)$ are $k(q)$-linearly independent. Suppose otherwise that there exist rational functions $a_k(q) \in k(q)$ such that
	\begin{align} \label{w:2}
		\sum_{i=1}^{k} a_k(q)\psi_q(v_i\otimes 1)=0\,. 
	\end{align}
	Without loss of generality we can assume that all $a_k(q)$ are non-zero. Let $\val_{q-1}$ be the valuation associated with $(q-1)$ and
	$$
	t=\min\{\val_{q-1}(a_1(q)), \dots, \val_{q-1}(a_k(q))\}\,.
	$$
	 Multiplying Equation \eqref{w:2} by $(q-1)^{-t}$ and applying $\res_{q=1}$ yields a non-trivial $k$-linear relation between elements $\res_{q=1}(\psi_q(v_i\otimes 1))$. The map $\psi_q$ restricts to  $\psi$, therefore $$\res_{q=1}(\psi_q(v_i\otimes 1))=\psi(v_i)\,.$$ This contradicts that $\psi$ is an  isomorphism.
\end{proof}

\end{document}
